\section{Testing coefficient relations by stable pairs}
\label{sec:stable-pairs}

We use the finite-dimensional correspondences of
Section~\ref{sec:correspondences} to prove
\begin{proposition}
\label{prop:kernel-inclusion}
For every datum and regular stability in Theorem~\ref{thm:main}, the
coefficient map induces a graded $R$-algebra surjection
\begin{equation}
\label{eq:stable-pair-surjection}
 B_F\twoheadrightarrow H_F^*(X).
\end{equation}
\end{proposition}
Throughout this section, $\mathsf K=G\times F$.  We retain the notation
$L$, $E_b$, $S_B$, $f_B$, $i=f_{\mathrm{pt}}$, and $\rho_Z$ from
Lemma~\ref{lem:correspondence-factorization}.  All cohomology in the test
below is ordinary equivariant cohomology.

\subsection{A degree constraint on global stable pairs}

We first record the geometric consequence of stability that will exclude
the first lattice of each word.  For a principal $G$-bundle $\mathcal P$
on $\mathbb P^1$, write $\mathcal L_\nu$ for its character line bundle,
with the convention that evaluating a $\nu^m$-semi-invariant on a section
of an associated representation gives a section of
$\mathcal L_\nu^{\otimes m}$.  Its degree is a linear function of the
topological class of $\mathcal P$.

\begin{lemma}
\label{lem:stable-section-degree}
Let $\mathcal P$ be a principal $G$-bundle on $\mathbb P^1$, equipped with
a frame at infinity, and let $\sigma$ be a global section of
$\mathcal P\times_G(T^*N)$.  Suppose that its value $\sigma_\infty$ in this
frame is $\nu$-stable and has trivial stabilizer.  Then
$\deg\mathcal L_\nu\geq0$.  If $\deg\mathcal L_\nu=0$, the pair
$(\mathcal P,\sigma)$ has a unique trivialization, agreeing with the given
frame at infinity, in which $\sigma$ is the constant section
$\sigma_\infty$.

In particular, suppose that $\mathcal P$ is specified by a lattice
$b\in\Gr_G$, with its usual trivialization on $\mathbb P^1\setminus\{0\}$.
In the degree-zero case its topological class is zero and
\begin{equation}
\label{eq:based-loop-trivialization}
 b=hG(\OO),\qquad
 h\in G(\C[t^{-1}]),\qquad h(\infty)=1.
\end{equation}
In the off-zero frame, the section is $h\sigma_\infty$.
\end{lemma}

\begin{proof}
Choose a $\nu^m$-semi-invariant $f$, with $m>0$, that is nonzero at
$\sigma_\infty$.
Its evaluation on $\sigma$ is a global section of a positive power of
$\mathcal L_\nu$, nonzero at infinity.  This proves nonnegativity of the
degree.

If the degree is zero, this section is nowhere zero.  Thus $\sigma$ is
everywhere semistable and determines a morphism
\[
 \gamma:\mathbb P^1\longrightarrow (T^*N)/\!/_{\nu}G.
\]
Every invariant polynomial evaluated on $\sigma$ is a regular function
on $\mathbb P^1$, hence constant.  The image of $\gamma$ therefore lies
in one fiber of the projective morphism to the affine quotient
$(T^*N)/\!/G$.  A sufficiently divisible GIT polarization is relatively
ample and pulls back to a positive power of $\mathcal L_\nu$.  Its
degree is zero, so $\gamma$ is constant: an ample line bundle has
positive degree on the image of a nonconstant morphism from
$\mathbb P^1$.  These are the usual character-linearized affine GIT
constructions; see \cite[Section~2]{King94}.

The constant value is the stable quotient point represented by
$\sigma_\infty$.  Its fiber in the semistable locus is the single orbit
$G\sigma_\infty$.  Consequently $\sigma$ is a section of the associated
orbit bundle $\mathcal P\times_G(G\sigma_\infty)$.  Since
$G\sigma_\infty\simeq G$ equivariantly, this section trivializes
$\mathcal P$.  Freeness makes the trivialization unique once the
coordinates of the section are prescribed to be $\sigma_\infty$; it
agrees with the specified frame at infinity.

In the lattice presentation, compare this global trivialization with the
fixed trivialization off zero.  Their transition is a regular map
$\mathbb A^1_{t^{-1}}\to G$ equal to $1$ at infinity.  It is the loop $h$
in \eqref{eq:based-loop-trivialization}, and the lattice of regular
sections at zero is $hG(\OO)$.
\end{proof}

Every point of $D=\mu^{-1}(0)^{\nu\text{-stable}}$ is stable with trivial
stabilizer also in the ambient representation $T^*N$.  Indeed, for a
closed invariant subvariety, semi-invariants lift by exactness of
reductive invariants, and orbit closures stay in that subvariety.
The affine GIT stability criterion therefore agrees with its ambient
restriction.  Lemma~\ref{lem:stable-section-degree} applies to the
\emph{pair} $(s,p)\in D$.

\subsection{The finite vanishing-cycle test}

For the truncation parameter $l$ of
Lemma~\ref{lem:correspondence-factorization}, put
\[
 L_-=\bigoplus_{j=-l}^{-1}t^jN,\qquad
 L_+=\bigoplus_{j=0}^{l-1}t^jN,\qquad
 P_-=\bigoplus_{j=-l}^{-1}t^jN^*.
\]
Thus $L=L_-\oplus L_+$, and the residue pairing identifies
$P_-$ with $L_+^*$.  Regard $s\in N$ and $p\in N^*$ as constant series.
On $\widetilde L=L\times D\times P_-$ define the $\mathsf K$-invariant
function
\begin{equation}
\label{eq:stable-test-function}
 W(n,s,p,y)=
 \Res\bigl(\langle p+y,n\rangle-\langle y,s\rangle\bigr),
 \qquad (s,p)\in D.
\end{equation}
After restricting $n$ to a moving lattice, variation in $y$ imposes the
prescribed nonnegative modes $\operatorname{pr}_{L_+}n=s$, while variation
in $n$ imposes dual-lattice membership for $p+y$.  These are the extension
conditions for the two components of the stable pair across the lattice
modification.
The residue is the coefficient of $t^{-1}$; the expression is
well-defined on the truncated spaces.

Use the shifted vanishing-cycle functor $\Phi_f=\phi_f[-1]$, normalized
so that $\Phi_0$ is the identity.  Vanishing cycles will be regarded
as complexes on the ambient space by closed direct image from the zero
fiber.  The test of a complex $\mathcal F$ on $L$ is
\begin{equation}
\label{eq:stable-test-functor}
 \mathsf T(\mathcal F)
 =R\Gamma_{\mathsf K}\!
 \left(\widetilde L,\,
 \Phi_W\,\operatorname{pr}_L^*\mathcal F\right).
\end{equation}
For a bounded projective lattice family $B$, set
$M_B=S_B\times D\times P_-$, and let $W_B$ be the pullback of $W$.
Proper base change and compatibility of vanishing cycles with proper
direct image give
\begin{equation}
\label{eq:stable-test-proper}
 \mathsf T(f_{B,*}\C_{S_B})
 \simeq R\Gamma_{\mathsf K}
       (M_B,\Phi_{W_B}\C_{M_B}).
\end{equation}
Here proper-direct-image compatibility is
\cite[Proposition~4.19(1)]{MaximSchurmann21}; the smooth projection
pullback uses part~(2).  The same operations in the critical-cohomology
setting are recalled in \cite[Appendix~A.1]{BottaTamagni26}.

The test is $R$-linear.  To see this precisely, work on a finite Borel
approximation with base $B\mathsf K$.  A coefficient class $u$ is a
morphism $\C_{B\mathsf K}\to\C_{B\mathsf K}[\deg u]$, pulled back to each
space.  The projection formula makes proper direct image compatible
with this action \cite[Theorem~6.3]{SchnurerSoergel16}.
Nearby cycles commute naturally with tensoring by a
locally constant complex pulled back from the base; taking the
specialization cone gives the same assertion for vanishing cycles.
Naturality with respect to the morphism representing $u$ therefore
identifies its tested action with cup product by the restricted class.
Passing to sufficiently large approximations proves the assertion for
$R=H_{\mathsf K}^*(\mathrm{pt})$.

\begin{lemma}
\label{lem:stable-test-support}
The support of $\Phi_{W_B}\C_{M_B}$ is contained in the closed locus
$\Sigma_B$ defined by
\begin{equation}
\label{eq:stable-critical-equations}
 n\in s+L_-,
 \qquad
 p+y\in gN^*(\OO)/t^lN^*(\OO),
 \qquad b=gG(\OO).
\end{equation}
Here $n\in E_b$ is already part of the definition of $M_B$.
If $B$ lies in a nonzero topological component whose character-line
degree is nonpositive, then $\Sigma_B$ is empty and the test in
\eqref{eq:stable-test-proper} is zero.
\end{lemma}

\begin{proof}
Differentiating in $y$ gives
$\operatorname{pr}_{L_+}n=s$.  Differentiating in the section variable
$n\in E_b$ says that $p+y$ annihilates $E_b$ in the residue pairing.
The annihilator is the displayed truncated dual lattice.  These
conditions are zeros of algebraic bundle maps and hence define a closed
locus.

This derivative calculation gives a support bound even when $B$ is
singular.  Off $\Sigma_B$, a derivative in one of the vector-bundle fiber
coordinates is nonzero.  The function is affine linear in that
coordinate; where its coefficient is invertible, replacing the
coordinate by $W_B$ gives a product chart over the unchanged base.
The constant sheaf has zero vanishing cycles in this chart.

A bounded lattice in each general linear factor specifies a modification
of the trivial vector bundle at zero, determined by finite jets.
A point of $\Sigma_B$ gives two global sections on $\mathbb P^1$ for the
resulting principal bundle.  In the off-zero frame they are the
polynomials $n$ and $p+y$ in $t^{-1}$, with values $s$ and $p$ at
infinity.  The lattice conditions are exactly their extendability at
zero: the discarded $t^l$ tails already lie in the corresponding
lattices.  Their value at infinity is the stable pair $(s,p)$.
Lemma~\ref{lem:stable-section-degree} now excludes negative
character-line degree and, in degree zero, forces the entire
topological class to be zero.  This proves the last assertion.
\end{proof}

\subsection{The neutral component and its quadratic normal form}

Let $Z$ be the neutral projective Schubert union used in
Lemma~\ref{lem:correspondence-factorization}, containing the base lattice
$1$.  Set $r=l\dim N$.  Since the index of the associated representation
lattice factors through $\pi_1(G)$, the bundles $E_b$ have rank $r$ on
$Z$.  Define the invariant open subset
\begin{equation}
\label{eq:stable-test-open}
 U=\{b\in Z:\operatorname{pr}_{L_+}:E_b\xrightarrow{\sim}L_+\}.
\end{equation}
It contains $1$.

\begin{lemma}
\label{lem:quadratic-test}
Assume $D\neq\varnothing$.  The locus $\Sigma_Z$ is a closed graph
$\Gamma\simeq U\times D$ in $M_Z$.  There are compatible
$\mathsf K$-equivariant identifications
\begin{equation}
\label{eq:quadratic-vanishing-cycles}
 \Phi_{W_Z}\C_{M_Z}
 \simeq \iota_{\Gamma,*}\C_\Gamma[-2r],
 \qquad
 \Phi_{W_{\mathrm{pt}}}\C_{M_{\mathrm{pt}}}
 \simeq \iota_{D,*}\C_D[-2r].
\end{equation}
Under these identifications, the test of $\rho_Z$ induces ordinary
restriction
\begin{equation}
\label{eq:quadratic-restriction}
 H_{\mathsf K}^*(U\times D)\longrightarrow H_{\mathsf K}^*(D)
\end{equation}
with the same uniform shift on source and target.  This map is split
surjective and $R$-linear.
\end{lemma}

\begin{proof}
First, every point of $\Sigma_Z$ lies over $U$.  By
Lemma~\ref{lem:stable-section-degree}, its principal bundle is trivial,
with a based loop $h$ as in \eqref{eq:based-loop-trivialization}.
A vector in the kernel of $E_b\to L_+$ is a negative polynomial section
of its associated $N$-bundle, extending at zero and vanishing at
infinity.  That associated bundle is trivial, so the vector is zero.
The two spaces have dimension $r$, and the projection is an
isomorphism.

Over $U$, write $E_b$ as the graph of a regular linear map
$A_b:L_+\to L_-$.  Thus $n=a+A_ba$ for $a\in L_+$.  Define the linear
functional
\[
 \ell_{b,p}(a)=\Res\langle p,A_ba\rangle.
\]
Via $P_-=L_+^*$, Equation~\eqref{eq:stable-test-function} becomes
\begin{equation}
\label{eq:stable-quadratic-coordinates}
 W_Z=\langle y,a-s\rangle+\ell_{b,p}(a).
\end{equation}
Its fiber critical equations have the unique solution
$a=s$, $y=-\ell_{b,p}$, depending regularly on $(b,s,p)$.
Consequently $\Sigma_Z$ is the graph $\Gamma\simeq U\times D$.
It is closed in all of $M_Z$ because the equations defining $\Sigma_Z$ are
closed and every solution has already been shown to lie over $U$.

It remains to compute the critical value.  Apply
Lemma~\ref{lem:stable-section-degree} to this solution.  In the
off-zero frame,
\[
 n=hs,\qquad p+y=hp.
\]
Invariance of the pairing gives
$\langle hp,hs\rangle=\langle p,s\rangle$, a constant series.  Hence
\[
 W_{\mathrm{crit}}
 =-\langle[t^{-1}](hp),s\rangle.
\]
Write $h=1+\xi t^{-1}+O(t^{-2})$, where $\xi\in\mathfrak g$.
The coefficient $[t^{-1}](hp)$ is the infinitesimal dual action
$\xi\cdot p$.  Therefore
\begin{equation}
\label{eq:stable-critical-value}
 W_{\mathrm{crit}}
 =-\langle\xi\cdot p,s\rangle
 =\langle p,\xi\cdot s\rangle
 =\langle\mu(s,p),\xi\rangle=0.
\end{equation}
Thus the regular, equivariant translations
$a'=a-s$, $y'=y+\ell_{b,p}$ turn
\eqref{eq:stable-quadratic-coordinates} into the fixed split form
\[
 W_Z=\langle y',a'\rangle
\]
over $U\times D$.

The vanishing cycles of this form on $L_+\oplus L_+^*$ are one copy of
$\C[-2r]$ at the origin.  The external-product theorem applies to the
constant sheaf of $U\times D$ even if $U$ is singular
\cite[Theorem, pp.~1--2 of the arXiv preprint]{Massey99}.  The splitting is global and
equivariant.  Its structure group acts as $g\oplus g^{-T}$, so the
quadratic vanishing-cycle line has trivial equivariant monodromy.
Choosing its identification with $\C$ once gives compatible
identifications over the whole base.

Lemma~\ref{lem:stable-test-support} says that the vanishing cycles are
zero off the closed graph $\Gamma$.  The product calculation therefore
gives the first isomorphism in
\eqref{eq:quadratic-vanishing-cycles} on all of $M_Z$.  In particular,
there are no further cycles over $Z\setminus U$: every point there
fails at least one of the fiber equations.  The same computation at
$1$ has $A_1=0$, $\ell_{1,p}=0$, and the same $2r$ quadratic variables,
giving the second isomorphism.

The map $\rho_Z$ is restriction to the base lattice.  To identify its
test as a morphism, not just its source and target, put $B=U\times D$
and consider the maps
\[
 D\xrightarrow{\,j\,}U\times D\xrightarrow{\,\pi\,}D,
 \qquad j(d)=(1,d),\quad \pi(b,d)=d.
\]
These maps are $\mathsf K$-equivariant and satisfy $\pi j=\id$.  Their ordinary
cohomological pullbacks satisfy $j^*\pi^*=\id$.
On the same base $B$, apply the natural external-product identification
to the morphism $\C_B\to Rj_*\C_D$, with zero function on $B$ and the
fixed split quadratic on $L_+\oplus L_+^*$.
Since $j$ is closed, proper-direct-image compatibility identifies the
target with the quadratic calculation over $D$.  One common choice of
generator of the quadratic line then identifies the resulting morphism
with $\C_B[-2r]\to Rj_*\C_D[-2r]$.
Taking ordinary equivariant cohomology gives exactly
Equation~\eqref{eq:quadratic-restriction}.  This uses naturality in a
sheaf morphism and proper direct image, not commutation with an arbitrary
nonsmooth pullback.  The identity $j^*\pi^*=\id$ supplies the
$R$-linear splitting, without any properness requirement on $\pi$.
\end{proof}

\begin{proof}[Proof of Proposition~\ref{prop:kernel-inclusion}]
The assertion is immediate if $D=\varnothing$, so assume $D$ is
nonempty.  Let $u\in R$ be homogeneous and zero in $B_F$.
Proposition~\ref{prop:minuscule} expresses $u$ in $A_F$ as a finite sum
of positive-length words in nontrivial dressed positive and negative
minuscule generators, with coefficients absorbed into the dressings.
Taking the full topological-degree-zero part, every word has total
topological degree zero.

The character-line degrees of its steps sum to zero.  Choose a step
with nonpositive character-line degree and, by commutativity, put it
first.  Its topological degree is nonzero.  Apply
Lemma~\ref{lem:correspondence-factorization}, with a common neutral
endpoint bound $Z$.  The morphism $u\rho_Z$ is a sum of compositions,
each factoring through the first-step object
$f_{B,*}\C_{S_B}$, with the prescribed shifts.  By
Lemma~\ref{lem:stable-test-support}, the test of every such object is
zero.  It follows that
\[
 \mathsf T(u\rho_Z)=0.
\]

The degree calculation in
Lemma~\ref{lem:correspondence-factorization} identifies the degree of
this morphism with $\deg u$, including when intermediate lattice
components are nonneutral.  The test commutes with shifts, and
Lemma~\ref{lem:quadratic-test} places the same quadratic shift on its
source and target.  Its $R$-linearity therefore gives the equality of
graded maps
\[
 (u|_D)\smile j^*=0.
\]
Apply this equality to $\pi^*(1)$.  Since $j^*\pi^*(1)=1$, we obtain
$u|_D=0$ in $H_{\mathsf K}^*(D)$.

The free stable-locus action identifies
$H_{\mathsf K}^*(D)=H_F^*(X)$, compatibly with the coefficient map.
We have proved
\[
 \ker(R\to B_F)\subseteq\ker(R\to H_F^*(X)).
\]
Both rings are generated by their coefficient images.  For $B_F$ this
follows from Proposition~\ref{prop:minuscule}; for $H_F^*(X)$ it follows
from Kirwan surjectivity \cite{MN18} and its equivariant extension
recalled above.  The displayed inclusion induces the graded $R$-algebra
surjection
\eqref{eq:stable-pair-surjection}.
\end{proof}
